\documentclass[11pt]{article}
\usepackage{mathblatt}
\begin{document}
% Kopfzeile Seite 1 ohne Kennung; die Rückseite (Lösungen, für den Lehrer) bekommt sie nach \begleitteil
\blattkopf{Mathematik}{Basisaufgaben}
\einheitenkopf[][Zettel 20]{Basisaufgaben · Zettel 20}
\zweigzeile{je Aufgabe 1 Punkt · Lösungen auf der Rückseite}
\par\vspace{2pt}\noindent Name \underline{\hspace{7cm}}\hfill Datum \underline{\hspace{3cm}}\hfill Punkte \underline{\hspace{1.2cm}} / 17\par\vspace{6pt}
% \small: volle Seite, kurze Aufgaben zweispaltig (v0.6)
\small

\noindent\begin{minipage}[t]{0.485\linewidth}\raggedright
% 1: Wert nach prozentualer Erhöhung berechnen – prozentrechnung-basis-k5-v3 (Original 2024-OS-B1e, ausgedacht: keine Marke)
\begin{aufgabe}{Eine Brezel kostet $2{,}50\,€$ und wird um $20\,\%$ teurer. Wie hoch ist der neue Preis? \leerfeld[€]}
\end{aufgabe}
\end{minipage}\hfill\begin{minipage}[t]{0.485\linewidth}\raggedright
% 2: Uhrzeit aus Startzeit und Dauer berechnen – einheiten-basis-k1-v3 (Original 2021-OS-B1a, ausgedacht: keine Marke)
\begin{aufgabe}{Ein Bus fährt um 10:27 Uhr ab, die Fahrt dauert $3$ h $48$ min. Wann kommt er an? \leerfeld Uhr}
\end{aufgabe}
\end{minipage}\par

\noindent\begin{minipage}[t]{0.485\linewidth}\raggedright
% 3: Größen vergleichen – einheiten-basis-k2-v6 (Original 2026-FOR-B1d, ausgedacht: keine Marke)
\begin{aufgabe}{Vergleiche $0{,}05$ m und $5$ cm. Setze das richtige Zeichen ein: $<$, $=$ oder $>$. $0{,}05$ m \leerfeld $5$ cm}
\end{aufgabe}
\end{minipage}\hfill\begin{minipage}[t]{0.485\linewidth}\raggedright
% 4: Zeiteinheiten umrechnen – einheiten-basis-k5-v6 (Original 2025-OS-B1b, ausgedacht: keine Marke)
\begin{aufgabe}{Ein Zeltlager dauert $2{,}5$ Tage. Gib die Dauer in Stunden an. \leerfeld[h]}
\end{aufgabe}
\end{minipage}\par

\noindent\begin{minipage}[t]{0.485\linewidth}\raggedright
% 5: Proportionale Zuordnung Dreisatz – zuordnungen-basis-k2-v6 (Original 2023-OS-B1a, ausgedacht: keine Marke)
\begin{aufgabe}{Ein Wasserhahn füllt $6$ l in $4$ min. Wie viele Liter sind es in $10$ min? \leerfeld[l]}
\end{aufgabe}
\end{minipage}\hfill\begin{minipage}[t]{0.485\linewidth}\raggedright
% 6: Lineare Gleichung lösen – lineare-gleichungen-basis-k1-v13 (Original 2024-OS-B1d, ausgedacht: keine Marke)
\begin{aufgabe}{Löse die Gleichung $2 \cdot (x - 7) = -4$. x = \leerfeld}
\end{aufgabe}
\end{minipage}\par

\noindent\begin{minipage}[t]{0.485\linewidth}\raggedright
% 7: Lösung durch Einsetzen prüfen – quadratische-gleichungen-basis-k1-v10 (Original 2025-OS-B1h, ausgedacht: keine Marke)
\begin{aufgabe}{Welcher Wert erfüllt $x \cdot (x - 4) = 12$? \\ \kreuz{$x = 2$} \kreuz{$x = 4$} \kreuz{$x = -2$} \kreuz{$x = -6$}}
\end{aufgabe}
\end{minipage}\hfill\begin{minipage}[t]{0.485\linewidth}\raggedright
% 8: Quadratseite aus Fläche berechnen – flaechen-basis-k2-v4 (Original 2020-OS-B1d, ausgedacht: keine Marke)
\begin{aufgabe}{Eine quadratische Tischdecke hat den Flächeninhalt $0{,}25$ m². Wie lang ist eine Seite? \leerfeld[m]}
\end{aufgabe}
\end{minipage}\par

\noindent\begin{minipage}[t]{0.485\linewidth}\raggedright
% 9: Umfang Rechteck berechnen – flaechen-basis-k5-v2 (Original 2026-FOR-B1h, ausgedacht: keine Marke)
\begin{aufgabe}{Ein Bilderrahmen ist rechteckig, $30$ cm lang und $20$ cm breit. Berechne den Umfang. u = \leerfeld[cm]}
\end{aufgabe}
\end{minipage}\hfill\begin{minipage}[t]{0.485\linewidth}\raggedright
% 10: Volumen Würfel berechnen – koerper-basis-k3-v3 (Original 2026-FOR-B1f, ausgedacht: keine Marke)
\begin{aufgabe}{Ein würfelförmiger Karton ist $10$ cm lang, breit und hoch. Gib sein Volumen an. \leerfeld[cm³]}
\end{aufgabe}
\end{minipage}\par

% 11: Vorzeichenregel anwenden – rationale-zahlen-basis-k2-v3 (Original 2015-OS-B1g, ausgedacht: keine Marke)
\begin{aufgabe}{Eine negative Zahl wird durch eine negative Zahl geteilt. Das Ergebnis ist … \\ \kreuz{immer positiv} \kreuz{immer negativ} \kreuz{Das kann man nicht entscheiden.}}
\end{aufgabe}

% 12: Wurzel eines Quadrats berechnen – potenzen-wurzeln-basis-k2-v3 (Original 2015-OS-B1j, ausgedacht: keine Marke)
\begin{aufgabe}{Welche Aussage ist wahr? Kreuze an.\\ \kreuz{$\sqrt{(-11)^2} = -11$}\\ \kreuz{$\sqrt{(-11)^2} = 11$}\\ \kreuz{$\sqrt{(-11)^2}$ ist nicht definiert}}
\end{aufgabe}

% 13: Prozent und Anteil umwandeln – prozentrechnung-basis-k2-v6 (Original 2022-OS-B1f, ausgedacht: keine Marke)
\begin{aufgabe}{Zwei von fünf Schülern fahren mit dem Bus – wie viel Prozent? Kreuze an. \\ \kreuz{$2\,\%$} \kreuz{$25\,\%$} \kreuz{$40\,\%$} \kreuz{$52\,\%$}}
\end{aufgabe}

% 14: Eigenschaft einer Figur zuordnen – winkel-dreiecke-basis-k5-v6 (Original 2026-FOR-B1c, ausgedacht: keine Marke)
\begin{aufgabe}{Kreuze die richtige Ergänzung an: „In jedem Parallelogramm …“\\ \kreuz{… sind die Diagonalen gleich lang.}\\ \kreuz{… sind alle Winkel gleich groß.}\\ \kreuz{… halbieren sich die Diagonalen gegenseitig.}}
\end{aufgabe}

% 15: Wahrscheinlichkeit einstufig – bruchrechnung-basis-k1-v9 (Original 2016-OS-B1f, ausgedacht: keine Marke)
\begin{aufgabe}{Ein Würfel mit acht Flächen trägt die Zahlen 1 bis 8. Wie groß ist die Wahrscheinlichkeit, eine Zahl größer als 5 zu werfen? \leerfeld}
\end{aufgabe}

% 16: Winkel über Scheitel- oder Nebenwinkel bestimmen – winkel-dreiecke-basis-k6-v1 (Original 2019-OS-B1a, ausgedacht: keine Marke)
\begin{aufgabe}{\begin{minipage}[t]{0.6\linewidth}Zwei Geraden schneiden sich. Wie groß ist $\varepsilon$? $\varepsilon =$ \leerfeld °\end{minipage}\hfill\begin{minipage}[t]{0.37\linewidth}\vspace{0pt}
\geradenkreuzung{40}{40^\circ}{}{\varepsilon}{}
\end{minipage}}
\end{aufgabe}

% 17: Term zu Figur angeben – flaechen-basis-k4-v3 (Original 2025-OS-B1i, ausgedacht: keine Marke)
\begin{aufgabe}{\begin{minipage}[t]{0.6\linewidth}Welche Formel beschreibt den Flächeninhalt $A$? \\ \kreuz{$A = p \cdot r : 2$} \kreuz{$A = (p + q) : 2$} \kreuz{$A = p \cdot q$} \kreuz{$A = p \cdot q : 2$}\end{minipage}\hfill\begin{minipage}[t]{0.37\linewidth}\vspace{0pt}
\dreieckrw{3.5}{2.5}{p}{q}{r}
\end{minipage}}
\end{aufgabe}

\begleitteil
\blattkopf{Basisaufgaben · Zettel 20}{BAS-Z20 · Lösungen}
\erg{1}{$2{,}50 \cdot 1{,}2 = 3{,}00\,€$ (nicht nur $0{,}50\,€$) \hfill {\footnotesize\textit{Wert nach prozentualer Erhöhung berechnen · 2024}}}
\erg{2}{10:27 Uhr $+ 3$ h $=$ 13:27 Uhr, $+ 48$ min $=$ 14:15 Uhr \hfill {\footnotesize\textit{Uhrzeit aus Startzeit und Dauer berechnen · 2021}}}
\erg{3}{$0{,}05$ m $= 5$ cm, also $0{,}05$ m $=$ $5$ cm \hfill {\footnotesize\textit{Größen vergleichen · 2026}}}
\erg{4}{$2{,}5 \cdot 24 = 60$ h \hfill {\footnotesize\textit{Zeiteinheiten umrechnen · 2025}}}
\erg{5}{Für $1$: $6 : 4 = 1{,}5$; $1{,}5 \cdot 10 = 15$ l \hfill {\footnotesize\textit{Proportionale Zuordnung Dreisatz · 2023}}}
\erg{6}{$x - 7 = -2$, also $x = 5$ \hfill {\footnotesize\textit{Lineare Gleichung lösen · 2024}}}
\erg{7}{$x = -2$, denn $(-2) \cdot (-2 - 4) = (-2) \cdot (-6) = 12$ \hfill {\footnotesize\textit{Lösung durch Einsetzen prüfen · 2025}}}
\erg{8}{$0{,}5 \cdot 0{,}5 = 0{,}25$, also $0{,}5$ m \hfill {\footnotesize\textit{Quadratseite aus Fläche berechnen · 2020}}}
\erg{9}{$u = 2 \cdot (30 + 20) = 100$ cm \hfill {\footnotesize\textit{Umfang Rechteck berechnen · 2026}}}
\erg{10}{$10 \cdot 10 \cdot 10 = 1\,000$ cm³ (nicht $10 \cdot 10 = 100$) \hfill {\footnotesize\textit{Volumen Würfel berechnen · 2026}}}
\erg{11}{immer positiv \hfill {\footnotesize\textit{Vorzeichenregel anwenden · 2015}}}
\erg{12}{$\sqrt{(-11)^2} = 11$, denn erst das Quadrat: $(-11)^2 = 121$, und die Wurzel ist nie negativ \hfill {\footnotesize\textit{Wurzel eines Quadrats berechnen · 2015}}}
\erg{13}{$40\,\%$ ($\frac{2}{5} = \frac{40}{100}$) \hfill {\footnotesize\textit{Prozent und Anteil umwandeln · 2022}}}
\erg{14}{… halbieren sich die Diagonalen gegenseitig. \hfill {\footnotesize\textit{Eigenschaft einer Figur zuordnen · 2026}}}
\erg{15}{$\frac{3}{8}$ (6, 7, 8) \hfill {\footnotesize\textit{Wahrscheinlichkeit einstufig · 2016}}}
\erg{16}{$\varepsilon = 40^\circ$ (Scheitelwinkel) \hfill {\footnotesize\textit{Winkel über Scheitel- oder Nebenwinkel bestimmen · 2019}}}
\erg{17}{$A = p \cdot q : 2$ \hfill {\footnotesize\textit{Term zu Figur angeben · 2025}}}
\end{document}
